\section{Minimum-defect normal forms and the hidden signed ternary torus}
\label{sec:bt3376-bt3389-cohomology-tau}
The flat $C_3^5$ port connections admit an exact minimum-defect word metric.  On one filled triangle there are $1$ states of length zero, $726$ states of length one, and $58{,}322$ states of length two, so the local enumerator is
\[
  1+726z+58{,}322z^2.
\]
Because the $240$ filled faces partition all $720$ edges, the global flat-space enumerator is its $240$th power and the exact flat-space diameter is $480$.  A sphere-counting argument against the $3^{2180}$ switching classes gives the certified quotient bound
\[
  \boxed{389\leq \operatorname{diam} H^1\leq480}.
\]
The upper endpoint is the universal minimum-defect normal form; the exact quotient covering radius inside this interval remains open.

The minimum defects also give a canonical $PSp(4,3)$-equivariant presentation.  For scalar coefficients,
\[
  0\longrightarrow \mathbb F_3[240\text{ faces}]
  \longrightarrow \mathbb F_3[720\text{ supports}]
  \longrightarrow Z^1\longrightarrow H^1\longrightarrow0,
\]
with dimensions $240,720,480,436$ after quotienting the $44$-dimensional coboundary space.  Tensoring by the externally labelled five-dimensional coefficient module gives $3600$ generators, $1200$ local face relations, $220$ coboundary relations, and
\[
  3600-1200-220=\boxed{2180}
\]
logical dimensions.  Thus the full logical sector is presented by minimum defects with exactly $1420$ independent relations.

For the affine involution
\[
 \tau(x_1,x_2,x_3,x_4,x_5)=(-x_4,1-x_3,1-x_2,-x_1,x_5)
\]
on $\mathbb F_3^5$, the $135$ orbit species have only $108$ distinct block-one-hot $Q_{15}$ barycentres.  Their fibre profile is
\[
  81\text{ singletons}+27\text{ double fibres}.
\]
The $27$ double fibres occur exactly at ternary displacement four.  Taking the coordinatewise missing symbols maps their barycentres bijectively to the $27$ fixed points of $\tau$.  The two orbit species above each such centre are distinguished by the invariant binary correlation
\[
 h=(x_1+x_4)(x_2+x_3-1)\in\{1,2\}.
\]
Consequently the $54$ ambiguous species form a canonical two-sheeted cover of the fixed affine flat $\mathbb F_3^3$.

More strongly, barycentric aggregation is exactly lumpable for the weighted Hamming-orbifold walk.  Its $135$-state operator splits as
\[
  \boxed{W_{135}\cong B_{108}\oplus D_{27}}.
\]
The barycentric block has spectrum
\[
10^1,\;7^6,\;4^{18},\;1^{32},\;(-2)^{33},\;(-5)^{18},
\]
and itself has the exact three-shell quotient
\[
 \begin{pmatrix}2&8&0\\2&4&4\\0&4&6\end{pmatrix}
\]
on shells of sizes $27,54,27$, with spectrum $10,4,-2$.

Orienting each double fibre by $h=1$ versus $h=2$ identifies the hidden block with the signed Cayley operator on $\mathbb F_3^3$
\[
  \boxed{D_{27}=C_3^{(3)}-C_3^{(1)}-C_3^{(2)}}.
\]
Equivalently it has weight $+1$ on the two $\pm e_3$ steps and weight $-1$ on the four $\pm e_1,\pm e_2$ steps.  Its exact spectrum is
\[
  4^4,\;1^{12},\;(-2)^9,\;(-5)^2.
\]
This is a finite signed-torus theorem, not a spacetime-signature claim.

Finally, a single minimum defect used as a $C_3^5$ voltage assignment generates only the subgroup $\langle v\rangle\cong C_3$.  The full $10{,}935$-vertex derived graph therefore splits into
\[
  \boxed{81\text{ connected components of }135\text{ vertices each}}.
\]
Each component is $32$-regular with $2{,}160$ edges and $15{,}714$ triangles.  Fourier decomposition across the $C_3$ fibre gives one untwisted $45$-state block and two conjugate magnetic blocks.  Their exact moments satisfy
\[
 \operatorname{tr}M^2=1440,\qquad \operatorname{tr}M^3=31{,}302,
\]
whereas the untwisted cubic trace is $31{,}680$.  The $378=42\cdot9$ deficit proves that this is a genuine nonseparable cohomological weighting.  It is a candidate chromatic-dual surface, not yet a certificate improving the bound nine.  Its cardinality match with the $135$ Hamming-orbit species is also not an isomorphism: the voltage component has degree $32$, while the Hamming-orbifold walk has degree $10$.
